\documentclass[11pt,article,oneside]{memoir}
\usepackage[utf8]{inputenc}
\usepackage[T1]{fontenc}
\usepackage{url}
\usepackage[hidelinks,colorlinks=true,linkcolor=blue]{hyperref}
\usepackage{amsmath,amssymb}
\usepackage{graphicx}
\usepackage{xcolor}
\usepackage{tabularx}
\usepackage{listings}
\usepackage{geometry}
\usepackage{titlesec}
\usepackage{libertine}

\geometry{margin=1in}

\titleformat{\section}{\Large\bfseries\color{darkgray}}{\thesection}{1em}{}
\titleformat{\subsection}{\large\bfseries\color{blue!80!black}}{\thesubsection}{1em}{}

\lstset{
  basicstyle=\small\ttfamily,
  breaklines=true,
  frame=lines,
  backgroundcolor=\color{gray!10}
}

\begin{document}

\begin{center}
    \Huge \textbf{Migration Guide} \\
    \vspace{0.5em}
    \large \textit{A CEREBRUM Research Publication (v1.2.4)} \\
    \vspace{1em}
    \large \textbf{Bryan Alexander Buchorn} \\
    Independent Researcher \\
    \vspace{2em}
\end{center}

\begin{abstract}
\textbf{Overview:} This guide provides a conceptual entry point into the CEREBRUM framework. It is designed for practitioners and researchers seeking to understand the core reasoning signals of Framework Implementation Guide.
\end{abstract}

\section{CEREBRUM Migration Guide}

This guide covers what operators and developers need to \textbf{do} (not just what changed) when upgrading between major versions. For a list of \textit{what} changed, see \texttt{CHANGELOG.md}.

---

\subsection{v1.0.0 $\to$ v1.1.0 (Phase 20: Relati\-vistic Hardening)}

\subsubsection{Required actions: none}
All four Phase 20 fixes use backward-compatible defaults. No existing code breaks.

\subsubsection{Recommended changes}

\textbf{1. Enable Query Snapshot Isolation (Hole 5)}

If your deployment uses the \texttt{GlobalRebalancer} for live graph updates alongside concurrent queries, enable snapshot isolation explicitly:

\begin{lstlisting}
# Before (v1.0.0 — no snapshot)
traversal = BeamTraversal(adapter=adapter, csa_engine=csa, ...)

# After (v1.1.0 — add snapshot at query start)
csa.set_query_snapshot(adapter.community_map)
traversal = BeamTraversal(adapter=adapter, csa_engine=csa, ...)
\end{lstlisting}

The \texttt{api/server.py} already does this autom\-atically for REST queries.

\textbf{2. Per-community CSA parameters for homogeneous graphs (Hole 6)}

If your graph has tightly-clustered domains (proteins, legal clauses, financial instruments) where intra-community similarity is uniformly high, add parameter overrides for those communities:

\begin{lstlisting}
# Identify homogeneous communities (avg intra-community S_C > 0.85)
# Then reduce β and increase α/γ for those communities:
csa = CSAEngine(
    community_params={3: (0.5, 0.15, 0.25, 0.05, 0.05, 0.0)},
    ...
)
\end{lstlisting}

\textbf{3. Path-preserving hold-out is now the default}

\texttt{InferenceValidator} now defaults to \texttt{path\_preserving=True}. On dense graphs this is a no-op. On sparse graphs (avg degree \textless  3) this will improve recall estimates. No code changes needed unless you speci\-fically want the old behavior:

\begin{lstlisting}
# To restore v1.0.0 behavior (not recommended):
validator = InferenceValidator(adapter, traversal, path_preserving=False)
\end{lstlisting}

---

\subsection{v0.4.0 $\to$ v1.0.0 (Phase 19: Production Hardening)}

\subsubsection{Required actions}

\textbf{1. Namespace isolation for mixed-modality deployments}

If you use both \texttt{IngestionPipeline} and \texttt{SignalEncoder} in the same graph, you must now add namespace prefixes to prevent entity collisions:

\begin{lstlisting}
# Before (v0.4.0 — collision risk)
pipeline = IngestionPipeline(adapter)
encoder  = SignalEncoder(entity_dim=128)

# After (v1.0.0 — isolated namespaces)
pipeline = IngestionPipeline(adapter, namespace="text")
encoder  = SignalEncoder(entity_dim=128, namespace="signal")
\end{lstlisting}

If your graph contains entities with the same name in both text and signal sources, they will now appear as separate nodes (\texttt{text:Name} and \texttt{signal:Name}). To inten\-tionally merge them, use \texttt{entity\_dedup\_map={"signal:Name": "text:Name"}} in \texttt{IngestionPipeline}.

\textbf{2. Wire the Zombie Bridge hook if using GlobalRebalancer + BridgeTwinEngine}

\begin{lstlisting}
# Before (v0.4.0 — stale bridges after rebalance)
rebalancer = GlobalRebalancer(adapter, q_drift_threshold=0.05)

# After (v1.0.0 — bridges stay consistent)
rebalancer = GlobalRebalancer(adapter, q_drift_threshold=0.05, bridge_engine=bridge_engine)
\end{lstlisting}

\subsubsection{Recommended changes}

\textbf{3. Enable warm-start for proba\-bilistic traversal on sparse graphs}

\begin{lstlisting}
# Before
traversal = BeamTraversal(adapter, csa, probabilistic=True)

# After — reduce cold-start variance 85%
traversal = BeamTraversal(adapter, csa, probabilistic=True, warm_start_strength=5.0)
\end{lstlisting}

\textbf{4. Enable causal flood protection for untrusted event streams}

\begin{lstlisting}
# Before
discretizer = STDPDiscretizer(w_threshold=0.5, n_min=5)

# After — adversarial hardening
discretizer = STDPDiscretizer(w_threshold=0.5, n_min=5, min_causal_span=1.0, use_chi_squared=True)
\end{lstlisting}

---

\subsection{v0.3.0 $\to$ v0.4.0 (Phase 18: v0.4 Horizon)}

\subsubsection{New required environment variable (if using API auth)}
\begin{lstlisting}
export CEREBRUM_JWT_SECRET="your-secret-key"
\end{lstlisting}

The API server now requires JWT authe\-ntication on all endpoints. Anonymous access can be re-enabled for development:
\begin{lstlisting}
export CEREBRUM_ALLOW_ANONYMOUS=true
\end{lstlisting}

\subsubsection{IngestionPipeline replaces direct adapter.add\_edge for text ingest}
\begin{lstlisting}
# Before (v0.3.0)
for subj, pred, obj in triples:
    adapter.add_edge(subj, obj, pred)

# After (v0.4.0) — preferred for text data
pipeline = IngestionPipeline(adapter)
for subj, pred, obj in triples:
    pipeline.process(subj, pred, obj, confidence=0.9)
\end{lstlisting}

Direct \texttt{add\_edge} still works — the pipeline is additive.

\subsubsection{LLM bridge import path changed}
\begin{lstlisting}
# Before
from cerebrum.bridge import generate

# After
from llm_bridge import generate, GenerationResult
\end{lstlisting}

---

\subsection{v0.2.0 $\to$ v0.3.0 (Phase 10–11: Production Hardening + Streaming)}

\subsubsection{igraph / leidenalg dependency removed}
The native Leiden reimp\-lementation (\texttt{core/leiden\_native.py}) is now the default backend. Remove these from your environment if installed:

\begin{lstlisting}
pip uninstall leidenalg python-igraph igraph
\end{lstlisting}

The native imple\-mentation is API-compatible; no code changes are needed.

\subsubsection{JWT authe\-ntication added to API}
See v0.3.0 $\to$ v0.4.0 section above (authe\-ntication was introduced in v0.3.0).

\subsubsection{StreamAdapter replaces manual polling loops}
\begin{lstlisting}
# Before (v0.2.0 — polling)
while True:
    event = get_next_event()
    adapter.add_edge(event.subject, event.object, event.predicate)
    time.sleep(0.1)

# After (v0.3.0 — StreamAdapter)
stream = StreamAdapter(base_adapter=adapter, window_size=10_000)
stream.push_event(subject, predicate, object, timestamp, weight)
\end{lstlisting}

---

\subsection{v0.1.0 $\to$ v0.2.0 (Phase 6–9: Federated)}

No breaking changes. All Phase 6–9 features (FederatedAdapter, Hologr\-aphicIndex, Handshake) are additive and do not affect single-graph deployments.

---
\textbf{Copyright © 2026 Bryan Alexander Buchorn. All Rights Reserved.}


\vspace{3em}
\begin{center}
    \small \textit{Project CEREBRUM — March 2026 — Hardened Enterprise Security (v1.2.4)}
\end{center}

\end{document}
